\section{过去信息的获取：观测、消息与记忆}
\label{sec:observer-theory}

固定系统边界和协议，设 $U\ne\varnothing$ 是同一装置所有允许历史的集合。目标比特 $Y(u)\in\{0,1\}$ 描述装置在历史 $u$ 中的过去。以 $H(u)$ 表示接收方在决策时可获取的信息，以 $W(u)$ 表示发送方掌握的信息，以 $Z(u)$ 表示实际送达的消息。只有当接收方仍能通过接口取得此前的读数，或已将其存入记忆时，这些读数才能计入 $H$。若两个协议允许访问的历史记录不同，它们定义的就是不同的问题。

\subsection{T1. 精确解码的判据}
\label{thm:observer-t1}

对任意观测映射 $O:U\to\mathcal O$，在所有允许历史上以确定性方式精确恢复 $Y$ 的充要条件为
\begin{equation}
 \forall u,v\in U\qquad O(u)=O(v)\ \Longrightarrow\ Y(u)=Y(v).
 \label{eq:observer-fibres}
\end{equation}
\emph{证明。} 若 $Y=d\circ O$，则对相同的输入应用同一函数 $d$，必得相同答案。反过来，对每个 $o\in O(U)$，集合 $\{Y(u):O(u)=o\}$ 非空，且由假设可知其中只有一个元素。令 $d(o)$ 等于该元素。$d$ 在 $O(U)$ 之外的值可任意规定。由此可得函数的存在性。在有限定义域上，只要存储容量和寻址能力足够，就能用查找表实现这一函数。

该判据可分别用于 $O=H$ 和 $O=(H,Z)$。要说明接收消息确实改善了零错误恢复能力，只须先找出 $H$ 相同而 $Y$ 不同的两段历史，再给出一个能由 $(H,Z)$ 对\emph{整个}历史类作精确解码的解码器。给定一族观测 $O_t$，能首次实现这种恢复的时刻，就是满足式~\eqref{eq:observer-fibres} 的最小 $t$，如果它存在。在具有擦除机制的系统中，判据在该时刻成立，仍不能确定它在更晚时刻是否成立。

\subsection{T2. 一条消息的代价}
\label{thm:observer-t2}

设 $U$ 有限。发送方通过无噪声信道发送一个定长码字 $c(W)$。发送与接收的时刻预先给定。以 $W(U)$ 为顶点集构造无向图 $G$。在顶点 $w,w'$ 之间连边的条件是存在 $u,v\in U$ 使
\[
 W(u)=w,\qquad W(v)=w',\qquad H(u)=H(v),\qquad Y(u)\ne Y(v).
\]
允许 $w=w'$。图中若有自环，任何这种形式的消息都不足以实现精确解码。若无自环，所需的最小字母表大小与最少传输比特数分别为
\[
 M_{\min}=\chi(G),\qquad b_{\min}=\lceil\log_2\chi(G)\rceil .
\]
\emph{证明。} 如果一条边的两个端点对应相同的消息，就会出现相同的 $(H,Z)$ 对应不同目标的情形，这与 T1 矛盾。因此，每种精确编码的取值都给出一个正常着色。自环会要求同一顶点具有两种不同的颜色。反过来，在任一正常着色下，$H$ 相同而目标不同的两段历史必定获得不同颜色。发送方只要传送颜色，便可由 T1 构造精确解码器。对 $M$ 种颜色作定长编码，最少且只需 $\lceil\log_2 M\rceil$ 个比特，包括 $M=1$ 时的空消息。

若发送方同时知道 $H$ 和 $Y$，就可以达到下界 $\max_{h\in H(U)}|\{Y(u):H(u)=h\}|$，因为它可以对每个 $h$ 分别给可能的目标编号。在对历史的访问受限时，该下界可能严格小于 $M_{\min}$。设接收方记录了有序对 $ab,ac,bc$ 之一。发送方知道从该有序对中曾经选中的那个字母，目标比特指明当时选中了第一还是第二个位置。每一对中的两个字母都需要不同的颜色。三个字母构成一个三角形，因此，尽管目标比特只有两个可能取值，这里仍需要三种消息符号。带有边信息的单次零错误编码与图着色之间的联系可追溯至 \href{https://doi.org/10.1109/TIT.1976.1055607}{Witsenhausen 的工作}。\href{https://doi.org/10.3390/e26040338}{Charpenay、Le Treust 和 Roumy} 还研究了利用边信息计算函数的问题。

\subsection{T3. 信道错误与带地址查询的代价}
\label{thm:observer-t3}

假设在历史 $u$ 中，受干扰的信号可以取非空集合 $\mathcal R(u)$ 中的任意值。对于\emph{每一种}允许的干扰都能实现恢复，当且仅当
\[
 H(u)=H(v),\quad Y(u)\ne Y(v)
 \quad\Longrightarrow\quad \mathcal R(u)\cap\mathcal R(v)=\varnothing .
\]
交集中的一个共同符号会使同一观测需要两个不同的答案。若上述交集均为空，则每个可达的 $(h,r)$ 对应唯一的目标值。这是在包含信道输出的扩展历史上应用 T1 的证明。

对于 $(H,R,Y)$ 的有限联合分布，若采用零一损失，最小错误概率衡量的是另一种性能，其值为
\[
 P_e^*=1-\sum_{h,r}\max_{y\in\{0,1\}}\Pr(H=h,R=r,Y=y).
\]
事实上，确定性答案 $d(h,r)$ 的成功概率为 $\sum_{h,r}\Pr(H=h,R=r,Y=d(h,r))$，每一项均可分别最大化。例如，设 $H,Z$ 为相互独立的均匀比特，$Y=H\oplus Z$，且 $R=Z\oplus N$，其中噪声 $N\sim\mathrm{Bernoulli}(p)$ 独立于前述比特，$0\le p\le 1/2$。最优答案 $H\oplus R$ 的错误概率为 $p$。只要 $p>0$，便无法对所有干扰保证精确恢复。

交互带来的优势取决于协议允许哪些消息。设 $A_0,A_1$ 为过去记录的两个不受限制的比特，接收方的本地地址 $K$ 决定所需的历史比特 $Y=A_K$。全部八种三元组 $(K,A_0,A_1)$ 都是允许的。一个查询比特 $K$ 和一个返回比特 $A_K$ 就已足够。假如返回的比特必须是与地址无关的固定函数 $c(A_0,A_1)$，那么仅凭该函数的一个值就能分别取 $K=0$ 和 $K=1$ 运行精确解码器，从而恢复整个有序对 $(A_0,A_1)$。四种不同的有序对至少需要四种消息。因此，在这一通信协议中，发送地址可以缩小返回消息所需的字母表。若允许返回两个比特，就可以不经查询直接传送完整记录。对于单个固定查询，若重放的正确答案与原答案具有相同的内容和送达时刻，确定性接收方将产生相同的结果。

\subsection{T4. 自身报告与保留状态}
\label{thm:observer-t4}

考虑两个装置，它们都精确地计算同一函数 $d(H,Z)=Y$ 并公布结果。第一个装置保存比特 $M=Y$。第二个装置擦除有关证据，使至少两段 $Y$ 不同的历史到达相同的可访问状态。随后，切断对先前公开报告以及观察者信道的访问。两个装置都接收一个新的比特 $C$，并被要求执行动作 $Y\oplus C$。第一个装置计算 $M\oplus C$。对于第二个装置，在上述两段历史中固定同一个 $C$。此时全部可获取的输入相同，所要求的比特却相反。由 T1 可知，第二个装置不存在精确解码器。因此，可以用延迟动作区分立即报告完全相同的两种实现。

在固定先前状态时对 $Z$ 施加干预，可检验动作是否依赖实际送达的信号。要检验计算结果能否跨轮次保留，还需要第二种干预：在新查询相同、且无法再访问先前报告的条件下，替换或擦除 $M$。这构成两种不同的反事实比较。在上述任务中，一个内部寄存器足以产生正确的延迟动作。
